\section{Discussion}\label{sec:discussion}

inferCNV is no longer actively supported or developed, and we do not plan further releases beyond the one described here. The repository README states this and points users to actively maintained alternatives, including CopyKAT \citep{gao2021copykat}, Numbat \citep{gao2023numbat} and InferCNA \citep{laffy_infercna}. We wrote this article so that the method in its final form is documented and can be cited, and so that results obtained with it can be interpreted with knowledge of what the code does.

inferCNV's main features have been: heatmaps of residual expression that make large-scale CNAs visible and easy to interpret; discrete state calls for up to more than two extra copies, calibrated by simulation to the data set at hand; confidence filtering of those calls; CNA-based subclustering; and a modular design in which smoothing, subclustering, HMM and filtering methods can be chosen and individual steps resumed. Its wide use across tumour types, platforms and, more recently, spatial assays suggests that this combination has been useful.

The method also has limitations that users should keep in mind. It requires a set of normal reference cells, chosen by the user, and its output is relative to them; if the reference cells themselves carry CNAs, or if the ``observations'' are mostly normal, the signal is distorted. It does not classify cells as malignant or normal by itself, and when nearly all cells are malignant the dominant tumour profile can be centred out as baseline unless additional normal cells are included \citep{song2025}. Its expression-based signal is affected by batch and platform effects \citep{chen2025}, and any expression difference between observations and references that spans many neighbouring genes can resemble a CNA. Schmid et al.~noted a higher sensitivity to gains than to losses \citep{schmid2025}. Per-cell calls are noisy, and analysis by subcluster, the default, is recommended. Finally, the Bayesian filtering step dominates run time on large data sets (Section 3.6). Allele-aware methods can resolve copy-neutral loss of heterozygosity and improve accuracy when allelic information is good \citep{gao2023numbat, chang2026preprint}; the release described here also includes an allele-aware extension of inferCNV by Yu Zhong, which will be described separately.
